\begin{wrapfigure}{r}{0.35\textwidth}
  \centering
  \vspace{-8pt}
  \begin{tikzpicture}[
      font=\sffamily\scriptsize,
      celE/.style={minimum size=3.6mm, inner sep=0pt, fill=serie-1, draw=papel, line width=\trazo, text=papel},
      celP/.style={minimum size=3.6mm, inner sep=0pt, fill=serie-3, draw=papel, line width=\trazo, text=tinta},
      titulo/.style={anchor=west, font=\sffamily\scriptsize\bfseries, text=uoc-azul},
      calle/.style={anchor=east, text=tinta-suave},
    ]
    % \bloque{y0}{filas}{etiquetas}: cuadrícula de tramos (E = entrenamiento, P = prueba), una fila por calle
    \newcommand{\bloque}[3]{%
      \foreach \s [count=\r from 0] in {#2} {%
        \foreach \l [count=\k from 0] in \s {%
          \node[cel\l] at (\k*3.6mm, #1-\r*3.6mm) {\l};}}%
      \foreach \t [count=\r from 0] in {#3} {\node[calle] at (-2.4mm, #1-\r*3.6mm) {\t};}}
    % Partición aleatoria: los tramos de una misma calle se reparten entre E y P
    \node[titulo] at (-14mm, 0) {Partición aleatoria};
    \bloque{-5mm}{{E,P,E,E,P},{E,E,P,E,E},{P,E,E,P,E},{E,E,E,P,E}}{Calle 1,Calle 2,Calle 3,Calle 4}
    % Validación agrupada: cada calle entera va a E o a P
    \node[titulo] at (-14mm, -22mm) {Validación agrupada};
    \bloque{-27mm}{{E,E,E,E,E},{P,P,P,P,P},{E,E,E,E,E},{P,P,P,P,P}}{Calle 1,Calle 2,Calle 3,Calle 4}
    % Validación entre ciudades: una ciudad entera para E, otra para P
    \node[titulo] at (-14mm, -44mm) {Validación entre ciudades};
    \bloque{-49mm}{{E,E,E,E,E},{E,E,E,E,E}}{Ciudad A,}
    \bloque{-58mm}{{P,P,P,P,P},{P,P,P,P,P}}{Ciudad B,}
    % Leyenda
    \node[celE, label={[anchor=west, text=tinta]right:Entrenamiento}] at (-12mm, -67mm) {E};
    \node[celP, label={[anchor=west, text=tinta]right:Prueba}] at (14mm, -67mm) {P};
  \end{tikzpicture}
  \caption{Reparto de los tramos entre entrenamiento y prueba según la estrategia.}
  \label{fig:tipos-particion}
  \vspace{-10pt}
\end{wrapfigure}
